\section{One Repository, Five Sittings}
\label{sec:walkthrough}

The repository here is a five thousand line Python service that answers price
quotes for a trading dashboard, and the first sitting is the hour
Figure~\ref{fig:two-records} draws. We assembled the sequence from sessions we ran
while building \sgt{} so that one repository contains the five operations
Section~\ref{sec:scenarios} asked for, in that order, followed by the one case
where \sgt{} stops and asks. Every command below runs and every block of output
came out of \sgt{} as it is today. Output lines too wide for this column are
wrapped, at the \ttdot{} separators where a line offers a list of alternatives,
and are otherwise unchanged. Read the sequence for which decisions the
developer is asked to make, and when in the work they arrive.

\subsection{Monday, 14:02 to 15:20: two saves and four names}

The developer made three requests in that hour, and the agent renamed a function
on its own initiative. At 14:33, with the rate limiting finished and the price
cache first drafted, the developer saved.

\begin{out}
\$ sgt save -m "rate limiting and a price cache on quotes"\\
\ttok{} save c39ab2f "rate limiting and a price cache on quotes"\\
~~\ttbranch{}\ttnew{} new feature (4a1c9e02) \ttdash{} unnamed; name it:\\
~~\ttvert{}~~~~sgt feature rename 4a1c9e02 "<label>"\\
~~\ttvert{}~~~~api.py::\_bucket, api.py::\_key, api.py::\_window +1 more\\
~~\ttbranch{}\ttnew{} new feature (001f73c2) \ttdash{} unnamed; name it: \dots\\
~~\ttvert{}~~~~cache.py::get\_cached, cache.py::\_ttl, db.py::query +1 more\\
~~\ttlast{}\ttnew{} new feature (9d20b18a) \ttdash{} unnamed; name it: \dots\\
~~~~~~~db.py::fetch\_row, api.py::export\_row, util.py::load\_row +1 more\\
~~\ttwarn{} one save touched 3 features \ttdash{} deliberate? `sgt log --map'\\
~~~~~shows them; `sgt feature regroup move' re-files work\\
~~\ttundo{} reverse this save:~~sgt undo
\end{out}

\noindent
The developer typed one sentence and \sgt{} wrote twelve records. The three groups
have no names yet, because \sgt{} writes records when the developer saves and works
out the names when the developer next reads, and a name costs a language model call
that has no business running while somebody is mid-session. The warning is the one
line here that resembles what git asks for, and \sgt{} printed it after the save was
already written. The developer read it, agreed the sitting had been mixed, and
carried on, and Section~\ref{sec:scenarios} opened with exactly that case.

The second save at 15:20 covered the retry policy and a rework of the cache the
agent made while adding the retry. By then the first save's groups had names, so its echo listed
\cmd{\ttdisc{} Price Cache (001f73c2)} for the group that already existed and
\cmd{\ttnew{} new feature (7c02b4d1)} for the retry work.
\cmd{sgt log --tree} afterwards printed the four groups the hour produced.

\begin{out}
Rate Limiting (4a1c9e02) \ttdot{} 4 symbol(s)\\
Price Cache (001f73c2) \ttdot{} 4 symbol(s)\\
Retry Policy (7c02b4d1) \ttdot{} 3 symbol(s)\\
Row Accessors (9d20b18a) \ttdot{} 5 symbol(s)
\end{out}

\noindent
Three of the four names are what the developer would have written. The fourth is
the agent's unrequested rename, named from the code because the developer never
said anything about that part of the work, the failure
Section~\ref{sec:names} predicts. It is also the only one of the four whose name
nobody needs, so nobody corrected it.

\subsection{Thursday: changing the rate limiter}

The fixed one minute window has to become a sliding one, and the developer starts
from what the rate limiting is now.

\begin{out}
\$ sgt show "Rate Limiting"\\
feature 4a1c9e02~~"Rate Limiting"\\
~~4 edits \ttdot{} 4 symbols in 1 file \ttdot{} last touched 3d ago\\
~~symbols~~~~~~api.py::\_bucket, api.py::\_key, api.py::\_window,\\
~~~~~~~~~~~~~~api.py::handle\\
~~saves~~~~~~~~c39ab2f~~rate limiting and a price cache on quotes\\
~~reverting this removes 4 edits
\end{out}

\noindent
The four functions are the answer to the question that took four runs of
\cmd{git log -L} in Section~\ref{sec:scenarios}, and nothing here required the
developer to guess that four was the right number of places to look. What this
output does not say is that the price cache calls \sym{api.py::\_bucket} and
therefore sits on top of the rate limiter, so a sliding window is a change
underneath work that already depends on it. The record holds that fact, because
the cache's edits name the functions they refer to, and \cmd{show} does not print
it. The developer finds it out on Wednesday instead, when \cmd{propose land}
refuses a subset for that reason, and a \cmd{show} that printed it here would
have cost us one line.

The agent rewrites \sym{api.py::\_window} and \sym{api.py::\_bucket} and leaves the
other two functions alone, and the developer saves, naming the group this time
because they already know which one it is.

\begin{out}
\$ sgt save --as "Rate Limiting" -m "sliding window, 60s of buckets"\\
\ttok{} save 1f0aee3 "sliding window, 60s of buckets"\\
~~\ttlast{}\ttdisc{} Rate Limiting (4a1c9e02)~~api.py::\_bucket, api.py::\_window\\
~~\ttok{} named "Rate Limiting"
\end{out}

\noindent
\cmd{--as} is the one place a developer states a grouping by hand, and it changed
nothing here, because both functions were already in that group and the edits would
have been filed there without it. They typed it because they knew the answer and
had no reason to let \sgt{} work it out. The rate limiting now has two stretches of
work on it, and \cmd{sgt log --focus 4a1c9e02} names both.

\begin{out}
\ttdisc{} Rate Limiting~~4a1c9e02~~\ttdot{}~~2 checkpoint(s)\\[0.35ex]
~~~[0\ttfull\ttfull\ttfull]~~4a1c9e02@0~~:rate-limiting~~Rate Limiting\\
~~~[1\ttfull\ttfull\ttempty]~~4a1c9e02@1~~:sliding-window~~Sliding Window\\[0.35ex]
~operate:\\
~~~~sgt revert 4a1c9e02:sliding-window~~(one checkpoint, by name)\\
~~~~sgt revert 4a1c9e02~~(whole feature)\\
~~~~sgt revert 4a1c9e02@1~~(by index)
\end{out}

\noindent
The developer never asked for these two stretches. \sgt{} cut the work where their
words changed, which in this repository is where one request ended and the next
began, and either half is addressable on its own, so
\cmd{sgt revert "Rate Limiting:sliding-window"} takes Thursday's work off and
leaves Monday's four edits in place. The bar beside each name puts that
checkpoint's edits in their place in the history, read left to right, with a middle
dot where a span of the history holds none of them, so the terminal spends four
characters on the encoding Figure~\ref{fig:two-records} spends a panel on. The last three lines exist because we watched
people on our own team discover that a stretch of work had a name and then stop,
not knowing which command would accept it.

\subsection{The following Monday: taking the cache out}

The cache serves a stale price for up to sixty seconds and has to go.

\begin{out}
\$ sgt revert "Price Cache"\\
~~~~[the log region redrawn, with the 6 removed edits marked]\\
~~subtracted from shared code (later work kept): api.py::handle\\
~~removed going forward: cache.py::get\_cached, cache.py::\_ttl,\\
~~~~~db.py::query\\
~~\ttwarn{} still references removed code (fix or revert separately):\\
~~~~~api.py::export\_row, util.py::load\_row, util.py::save\_row\\[0.35ex]
apply this revert? [y/N]
\end{out}

\noindent
The retry policy edited \sym{api.py::handle} after the cache had edited it, so
that function keeps the retry work and loses the cache's three lines, and a
request made a week ago comes out from under a change made yesterday. Two
functions existed only for the cache and come out with it, which empties
\sym{cache.py} and removes the file, and
\sym{db.py::query} loses the parameter the cache added and goes back to the text it
had at 14:02. The three functions in the last line are row accessors the agent
rewrote during its rename, and it pointed all three at \sym{get\_cached} while it
was in there, so their edits belong to a request the developer is keeping and name
a function the revert deletes. \sgt{} names them and proceeds, because whether the
repair is deleting three call sites or reverting the rename as well is not a
question about the recorded edits.

The developer answered \cmd{y} and ran the test suite. It failed in the three
places the report named and nowhere else, and the repair was deleting the cache
lookup from each of them. \sgt{} builds every one of these operations as set
arithmetic so that a revert can be taken back exactly, and the developer tested
that claim before trusting it.

\begin{out}
\$ sgt restore "Price Cache"\\
~~\ttok{} restore applied \ttdash{} 0 edit(s) removed, 6 added.\\
~~~~(`sgt undo' reverses this.)
\end{out}

\noindent
Putting the cache back returned \sym{cache.py}, \sym{db.py::query}, and
\sym{api.py::handle} to the bytes they held before the revert, and it left the
three call site repairs the developer had made in between alone, because a restore
adds records to the set and rebuilds the files from the result without reapplying
a patch. They had wanted the cache gone, so they reverted it a second time, and
that run printed the same three consequence lines as the first.

\subsection{Wednesday: shipping the export fix alone}

A bulk export endpoint has to ship to the release branch today, and the branch it
sits on also carries the retry policy, which is half finished. The export code goes
into a new \sym{export.py} and applies the rate limit before streaming, so
\sym{export.py::write\_csv} calls \sym{api.py::\_bucket}. Nobody thought about that
call as a dependency while writing it.

\begin{out}
\$ sgt propose create --base release\\
\ttok{} proposal 3f9a2c81b04d (base release, +17 op(s), 4 feature(s))\\
\$ sgt propose land 3f9a2c81b04d --subset "Quote Export"\\
\textnormal{$\times$} 'Quote Export' requires Rate Limiting -- include it in\\
~~--subset too
\end{out}

\noindent
The developer asked to move one group and \sgt{} named the group that has to come
with it, the fourth operation Section~\ref{sec:scenarios} asked for. The record of the call to \sym{api.py::\_bucket} was
written at the moment the export code was saved, and \sgt{} reads it now to answer
a question nobody was thinking about then. Git would have taken the cherry-pick and
the release branch would have failed to import.

\begin{out}
\$ sgt propose land 3f9a2c81b04d --subset "Rate Limiting" "Quote Export"\\
\ttok{} propose land release: 8b21c0d4e7fa (+8 op(s))
\end{out}

\noindent
Two groups went to the release branch in the order the records put them in, and the
half finished retry policy and the accessor rename stayed behind. Those two groups
are the stack from Section~\ref{sec:scenarios}, assembled without the branch that
had to exist before the work did.

\subsection{Thursday: two versions of one function}

A colleague on this repository had been editing \sym{api.py::handle} as well,
adding a request identifier to every outbound log line. The developer syncs.

\begin{out}
\$ sgt sync\\
\ttwarn{} sync origin/main: merged 4f2a1c8b93de\\
~~~~+7 op(s)\\
~~~~\ttwarn{} 1 OPEN FORK(S) \ttdash{} forked symbol(s) sit at the common\\
~~~~~~ancestor until resolved:\\
~~~~~~api.py::handle \ttarrow{} sgt resolve api.py::handle
\end{out}

\noindent
The remedy is printed as a function name because that is what the developer knows
about the situation. Underneath it there is a command that takes the two edge
identifiers directly, \cmd{sgt advanced merge-op 3c91ab02 8ee14d77}, and for the
first several months of this system's life that was the line the fork warning
printed, on the reasoning that it named the two things precisely. A developer who
reads it learns that two records are in conflict and still has to work out which
function they are in. Deriving the function's name and printing that instead is a
one-line change we made late, and it is the clearest instance in this repository of
a report that was accurate and useless. The two
versions of \sym{api.py::handle} are the one the cache revert produced, three lines
shorter at the top where the cache lookup had been, and the colleague's, one line
longer inside the logging block at the end. No line is common to the two changes,
so git merges them without comment, the case
Figure~\ref{fig:fork} draws. \sgt{} prints them beside each other and asks.

\begin{out}
\ttrule{} fork on api.py::handle \ttrule{}~~(\ttleft{} left = tip A~~\ttright{} right = tip B)\\
~~~~[the two versions of the function, side by side]\\
pick [l]eft / [r]ight / [e]dit / [q]uit:
\end{out}

\noindent
The developer picked \cmd{e}, merged the two by hand in about thirty seconds, and
\sgt{} ran the project's test command before including the result. We are not going
to claim the textual merge would have been wrong here. It would most likely have
been right, and \sgt{} charged thirty seconds and one interruption to establish
that. We designed for the case where the merge is wrong, one edit adding an early
return and the other edit's new code sitting below it, and the test run that would
have caught it happening on a branch nobody has pushed yet. A developer
who merges often and reads carefully is paying us for a certainty they already had.

\subsection{What this cost}

Across five sittings and one fork the developer made three decisions \sgt{} could
not make for them. They decided the export endpoint should ship without the retry
policy, they decided how the two versions of \sym{api.py::handle} should read, and
they decided that the three functions still calling the cache should lose that
call and keep the rest of the agent's rename. All three are decisions about the
code itself, and none of the three is a question about which lines belong in which
commit.

The developer paid for those three decisions in three places. One group of five
edits carries a name derived from the code, because nobody said anything about
that part of the work. The revert left \sym{api.py::handle} holding text nobody
had read, so their test suite decided whether the revert had worked. The fork
stopped them on a merge that was probably fine. The second and third of those are direct consequences of
Section~\ref{sec:sets} and Section~\ref{sec:fork}, and on a repository this size we
would choose both again. On a repository where twenty people edit the same
functions each day, the fork rule is the first thing we would expect to have to
change.
